\documentclass[11pt,a4paper]{article}
\usepackage[utf8]{inputenc}
\usepackage{amsmath,amssymb}
\usepackage{graphicx}
\usepackage{booktabs}
\usepackage{geometry}
\geometry{left=2.5cm,right=2.5cm,top=2.5cm,bottom=2.5cm}
\usepackage{natbib}

\title{A likelihood ratio test for q-Weibull model selection under interval censoring}

\author{Huirong Cao$^{1}$, Fuchang Wang$^{2}$\\
$^{1}$ College of Science, Langfang Normal University, Langfang, China\\
$^{2}$ Emergency Management University, Beijing, China}

\date{}

\begin{document}
\maketitle

\begin{abstract}
The q-Weibull distribution extends the Weibull by introducing a shape deformation parameter $q$ that controls tail heaviness ($q>1$) or bounded support ($q<1$). Under interval censoring, estimating $q$ is challenging because the tail information is partially observed. We propose a likelihood ratio test for comparing q-Weibull against Weibull, and show via simulation that the test has near-nominal size and reasonable power. We apply the test to a large dental emergence study (n=4386) and find significant evidence against the Weibull in favor of the q-Weibull with $q>1$, consistent with unobserved heterogeneity.
\end{abstract}

\section{Introduction}

The Weibull distribution is the standard parametric model in survival analysis and reliability engineering. Its survival function is $S(t)=\exp(-\lambda t^\beta)$. However, the Weibull has a fixed exponential tail and unbounded support. Real lifetime data often exhibit either heavier tails (due to unobserved heterogeneity) or bounded maximum lifetimes.

The q-Weibull distribution \citep{picoli2003q} introduces a deformation parameter $q$:
\begin{equation}
S_q(t) = \left[1 + (q-1)\lambda t^\beta\right]^{-1/(q-1)}, \quad q \neq 1,
\end{equation}
which reduces to the Weibull as $q \to 1$. When $q>1$, the tail decays as a power law (heavier than Weibull); when $q<1$, the distribution has a finite upper endpoint $T_{\max} = [1/((1-q)\lambda)]^{1/\beta}$.

Interval-censored data arise when the exact event time is unknown, but it is known to lie in an interval $(L_i, R_i]$. This is common in periodic follow-up studies. Under interval censoring, the tail of the distribution is less precisely identified, making the estimation of $q$ particularly difficult.

We propose a likelihood ratio test for $H_0: q=1$ (Weibull) vs $H_1: q \neq 1$ (q-Weibull), under interval censoring. We evaluate the test via simulation and apply it to a large dental dataset.

\section{Methods}

\subsection{Model}

For interval-censored observations $(L_i, R_i]$, the likelihood contribution is $S(L_i) - S(R_i)$. For right-censored observations at $L_i$, the contribution is $S(L_i)$.

We fit both the Weibull ($q=1$) and q-Weibull (free $q$) by maximum likelihood. The likelihood ratio statistic is:
\begin{equation}
\Lambda = 2 \left[\ell(\hat{q}, \hat{\lambda}, \hat{\beta}) - \ell(q=1, \tilde{\lambda}, \tilde{\beta})\right].
\end{equation}

Under $H_0$, $\Lambda$ is approximately $\chi^2_1$. However, because interval censoring reduces tail information, the asymptotic approximation may be poor in small samples. We also provide a parametric bootstrap alternative.

\subsection{Interpretation}

- If $q < 1$ is supported: the data suggest bounded support.
- If $q > 1$ is supported: the data suggest a heavier-than-Weibull tail. Note that this can arise from either genuine tail deformation or unobserved heterogeneity (frailty), which produces the same survival function form \citep{hougaard2000frailty}.

\section{Simulation study}

We generate interval-censored data mimicking annual examinations: $L = \lfloor t \rfloor$, $R = L+1$, with 37\% right censoring. We set $\lambda=1/16$, $\beta=2$, and vary $q$. We run 200 replications per scenario and use the asymptotic $\chi^2_1$ reference distribution.

\begin{table}[h]
\centering
\caption{Simulation results (200 reps). Size under Weibull null.}
\begin{tabular}{lcccc}
\hline
Scenario & $\hat{q}$ (mean$\pm$sd) & $\Lambda$ (mean) & Rejection rate ($\alpha=0.05$) \\
\hline
Weibull, n=200 & 0.93 $\pm$ 0.28 & 1.13 & 0.070 \\
q=0.8, n=200 & 0.56 $\pm$ 0.29 & 2.59 & 0.255 \\
q=0.7, n=200 & 0.46 $\pm$ 0.25 & 3.57 & 0.365 \\
q=0.6, n=200 & 0.34 $\pm$ 0.21 & 4.84 & 0.580 \\
Weibull, n=500 & 0.94 $\pm$ 0.16 & 0.99 & 0.030 \\
q=0.7, n=500 & 0.41 $\pm$ 0.19 & 8.70 & 0.840 \\
\hline
\end{tabular}
\end{table}

The size is close to the nominal 5\% (7.0\% at n=200, 3.0\% at n=500). Power increases with sample size and with deviation from $q=1$. At n=500 and q=0.7, power reaches 84\%.

\section{Real data application}

We apply the test to the Signal Tandmobiel dental emergence study (n=4386), where children were examined annually and tooth emergence times are interval-censored.

\begin{table}[h]
\centering
\caption{Real data results.}
\begin{tabular}{lcc}
\hline
Model & nll & Parameters \\
\hline
Weibull & 5597.0 & $\lambda$, $\beta$ \\
q-Weibull & 5576.7 & $q$, $\lambda$, $\beta$ \\
\hline
LR statistic & \multicolumn{2}{c}{$\Lambda = 40.6$, $p < 0.0001$} \\
Estimated $q$ & \multicolumn{2}{c}{$\hat{q} = 1.35$} \\
\hline
\end{tabular}
\end{table}

The test strongly rejects the Weibull in favor of the q-Weibull. The estimated $q=1.35$ indicates a heavier-than-Weibull tail. As noted above, this is consistent with either genuine tail deformation or unobserved heterogeneity (frailty). The q-Weibull survival function with $q>1$ is mathematically equivalent to a gamma frailty model with $\theta = 1/(q-1) \approx 2.86$.

\section{Discussion}

We have proposed a likelihood ratio test for q-Weibull model selection under interval censoring. The test has near-nominal size and reasonable power in moderate samples. The real data application demonstrates that the test can detect deviations from the Weibull assumption.

However, two caveats are important. First, a significant test in the direction $q>1$ does not distinguish between tail deformation and frailty. Second, when $q<1$, the test detects bounded support, but estimating the exact upper endpoint remains difficult under interval censoring.

Future work will develop a formal test for distinguishing $q>1$ tail deformation from frailty, perhaps using additional model features such as the hazard function shape.

\bibliographystyle{apalike}
\begin{thebibliography}{99}

\bibitem[Picoli et al.(2003)]{picoli2003q}
Picoli, I., Mendes, R.S., and Malacarne, L.C. (2003). q-Weibull distributions: generalized Weibull distribution. \emph{Physica A}, 327(1-2), 151--158.

\bibitem[Hougaard(2000)]{hougaard2000frailty}
Hougaard, P. (2000). \emph{Analysis of Multivariate Survival Data}. Springer.

\end{thebibliography}

\end{document}
